\chapter{INTRODUCTION}

{\baselineskip=2\baselineskip
Short introduction ...


\section{Background of the Study}

Advanced technologies have become a cornerstone of modern agriculture, streamlining production and quality control through data-driven methods. Recent trends emphasize the integration of machine learning to enhance efficiency through precision farming and improved resource allocation (Botero-Valencia et al., 2025). Research indicates a paradigm shift where machine learning models have been useful in recent systems dominating over food production. These approaches allow for the evaluation of internal fruit characteristics without compromising physical integrity, aligning practices with contemporary food safety and environmental management requirements (Araújo et al., 2023).

The shift toward quantitative analysis is further supported by research into automated grading systems. Studies highlight how machine learning algorithms are transforming traditional agricultural practices by providing high-precision sorting and quality assurance (Botero-Valencia et al., 2025). Furthermore, the integration of computer vision and acoustic analysis has proven effective in handling agricultural products, as demonstrated utilized specialized algorithms to provide accurate and efficient solutions for produce management  (V. et al., 2025). These trends illustrate a growing reliance on feature extraction techniques to interpret complex biological data.
In 2023, the Philippines ranked \#52 in the global production of watermelon, having produced around 124,737.55 metric tons of watermelon (Trends of Fresh Watermelon Production in Philippines, 2024). Although the Philippines’ contribution to the total global production of watermelon is only around 0.12\%, it is still a high-value crop in the Philippines. 

Despite these global advancements, many local farmers in the Philippines still rely heavily on traditional, manual methods for pre-harvest assessment. Currently, determining the maturity of watermelons in the field depends on the thumping technique—manually tapping the fruit and listening to the resulting sound, in addition to visual inspection and the observation of the wilting of the watermelon’s tendrils. This practice is subjective, as the interpretation of acoustic resonance varies significantly among individuals. Recent literature notes that subjective harvest decisions can lead to the collection of underripe produce, which lacks the necessary sugar content and flavor profiles required for market standards (Savaş Koç \& Ferit Akbalik, 2025). This reliance on human intuition creates an inconsistency in fruit quality that can affect the overall reliability of the harvest. 

To address these challenges, this study proposes a portable system for detecting watermelon ripeness using acoustic signal processing and machine learning. The device will utilize a microphone to capture the acoustic signals produced by tapping the fruit while still on the vine, with a camera for cross-validation.

\section{Statement of the Problem}

This study seeks to investigate the challenges associated with accurate and objective watermelon ripeness assessment. Despite significant advances in agricultural sensing and automation, accurate and objective watermelon ripeness assessment remains a persistent challenge, particularly for smallholder farmers who continue to rely on manual and subjective methods. These experience-based practices lack standardized metrics and fail to reliably predict quality attributes such as soluble solids, sugar, acidity, sweetness, and firmness (Shah et al., 2025). 

To overcome these inconsistencies, various nondestructive testing techniques have been explored, including acoustic-based methods such as acoustic/tactile cues (Hamid et al., 2023), FFT-based acoustic tapping (Okafor, 2020), acoustic impulse response (Pamungkas & Bintoro, 2021), and frequency-domain analysis combined with machine learning (Li et al., 2025), as well as computer vision-based methods such as color, texture, and shape feature extraction (Jie et al., 2024), YOLO-based deep learning detection (Defang et al., 2024), and UAV-mounted vision systems (Singh et al., 2025).

Among these, acoustic-based methods have gained particular attention due to their low cost, portability, and ease of field implementation, yet their performance remains highly sensitive to environmental noise with wind, machinery, and handling interference significantly degrading signal quality and the reliability of time-domain and frequency-domain features used for ripeness classification (Li et al., 2025). Similarly, computer vision and optical-based methods, while effective at detecting external ripeness indicators such as color, texture, and shape, are fundamentally limited to surface-level assessment and cannot directly measure internal quality attributes such as flesh firmness, sugar content, or internal voids (Jie et al., 2024; Defang et al., 2024; Arboleda et al., 2020). Relying on either approach alone is therefore insufficient for accurate and comprehensive ripeness assessment.

Existing approaches address either acoustic signal processing or computer vision-based analysis in isolation, with no portable, integrated system combining both modalities for real-time field use identified in the literature. Consequently, local farmers continue to face post-harvest losses and inconsistent product quality due to poor handling, lack of technology, and inadequate ripeness assessment tools (Adepoju & Ologan, 2020). 

This study addresses this gap by developing a portable device that analyzes the acoustic response and tendril visual characteristics of watermelons, enabling objective and real-time ripeness classification directly in the field and providing local farmers with a practical and reliable alternative to subjective manual assessment. 

\section{Objectives of the Study}

In view of the above stated problem, we have the following objectives:
\begin{enumerate}
	\item To design a system using an acoustic sensing module utilizing digital signal processing to capture vibration signals and isolate resonant frequencies and a camera to capture tendril images;
	\item To develop a portable, field-deployable unit that will analyze acoustic and visual inputs for immediate on-site classification;
	\item To evaluate the system’s classification accuracy and reliability by comparing its output against traditional manual methods and destructive quality validation. 

\end{enumerate}

\section{Significance of the Study}

The development of a portable, non-destructive ripeness assessment system provides critical advancements in agricultural technology and quality control. The findings and output of this research will be significantly beneficial to the following: 

\textbf{Local Farmers and Agricultural Producers.} This study provides a low-cost, objective tool that replaces subjective "thumping" methods. By providing accurate harvest timing, farmers in regions such as Northern Mindanao can minimize post-harvest losses, ensure consistent fruit quality, and command better market prices in urban centers like Cagayan de Oro City.

\textbf{Consumers and Retailers.} The system ensures that only watermelons at peak maturity reach the market. This increases consumer confidence and satisfaction by guaranteeing internal quality attributes—such as sweetness and firmness—reducing the financial risk of purchasing over-ripe or hollow-heart produce.

\textbf{Future Researchers.} The data and methodologies established in this study can serve as a baseline for further research into other fruit varieties or the integration of more complex artificial intelligence models for real-time agricultural diagnostics in uncontrolled environments.



\section{Scope and Limitations}

The study will be conducted in a specific watermelon farm in Northern Mindanao as the site for data collection and system testing. The proposed system will be designed to function as a portable, handheld tool capable of evaluating watermelon ripeness prior to harvest. It relies on a combination of acoustic signal processing and image-based analysis, allowing it to evaluate both external and internal indicators of maturity. All data used for system development and testing will be gathered under natural field conditions to ensure that the device reflects how it would function in actual farm uses.

While the system is designed to provide reliable ripeness assessment, its performance can still be influenced by environmental factors such as wind, machinery, or human activity, which may affect acoustic readings. In addition, the device evaluates only one melon at a time, limiting its use for simultaneous or large-scale assessments. Its focus is exclusively on pre-harvest ripeness, leaving post-harvest quality, storage behavior, and shelf life beyond its scope. Moreover, performance is shaped by the specific farm conditions and the watermelon variety used during development, meaning that results may differ when applied to other environments or cultivars.


\section{Definition of Terms}

\begin{description}

	\item[Acoustic Signal Processing.] 
	A technique used to capture and analyze the sound waves generated from a watermelon "thump" to identify frequency patterns associated with maturity.
	
	\item[Machine Learning.] 
    A category of computational algorithms used to classify watermelon ripeness by identifying patterns within the extracted features.

    \item[Pre-harvest Assessment.]
    The process of evaluating the maturity of watermelons while they are still attached to the vine to determine the optimal timing for harvest.

    \item[Non-Destructive Testing.] 
    A method of quality evaluation that allows for the inspection of the internal state of a fruit without causing physical damage.

    \item[Thumping Technique.]
    The manual act of striking a watermelon with the hand to produce the resonant sound required for microphone input and digital analysis.
	

\end{description}

}
